\begin{lemma}
\label{thm: fixed points are on the boundary}
    If the learning rates $(\lrate_k)$ and the update maps $(\updateset_k)$ satisfy \autoref{asp: stochastic conditions on effective learning rate}, 
    a sequence $(q_k)$ generated by MC-O-PI converges with probability one to a suboptimal fixed point 
    if and only if that point lies on the boundary between several policy regions.
\end{lemma}
\begin{proof}
    Let the sequences $(q_k)$ and $(\pol_k)$ be generated with MC-O-PI, 
    and suppose $(q_k)$ converges with probability one to a suboptimal fixed point $\hat{q}$. 
    If $\hat{q}$ is not on the boundary between several regions, there exists a unique policy $\hat{\pol}$ and an $\epsilon > 0$ such that 
    \begin{align*}
        \forall q \in \qset : 
        \quad
        \|\hat{q} - q\| \leq \epsilon
        \implies
        \avalset(q) = \{q_{\hat{\pol}}\} .
        % \forall s \in \sset
        % \hat{\pol} = \argmax_{\pol \in \polset} \hat{q}(s, \pol) = \argmax_{\pol \in \polset} q(s, \pol) .
    \end{align*}
    Furthermore, the definition of convergence implies a time step $K$ such that for all $k \geq K$
    \begin{align*}
        \|\hat{q} - q_k\| \leq \epsilon
    \end{align*}
    with probability one.
    As a result, for all $k \geq K$ 
    % it also holds that $q_{\pol_k} = q_{\hat{\pol}}$ with probability one.
    system \eqref{eq: complete difference inclusion} simplifies to
    \begin{align}
        \label{eq: mcopi for polk}
        q_\kpo \in \big\{ q_k + \lrate_k \update_k (q_{\hat{\pol}} - q_k + w_k) : \update_k \in \updateset_k(q_k) \big\} .
    \end{align}
    Proposition~4.1 and Example~4.3 in \citep{Bertsekas1996Neuro-DynamicProgramming}  guarantee that, under \autoref{asp: stochastic conditions on effective learning rate}, every sequence generated by system \eqref{eq: mcopi for polk} converges to $q_{\hat{\pol}}$ with probability one, which implies that $\hat{q} = q_{\hat{\pol}}$.
    However, as noted in \autoref{rmk: only qopt is inside its own region}, if $\hat{q}$, and thus $q_{\hat{\pol}}$, are suboptimal, then $\hat{q} \notin \region(\hat{\pol})$, which contradicts the initial assumption.
    As a result, if the sequence $(q_k)$ converges with probability one to a suboptimal fixed point, it must lie on the boundary between several regions.
\end{proof}
